\documentclass[10pt,a4paper]{amsart}
\usepackage[margin=19mm]{geometry}
\usepackage{amsmath,amssymb,mathtools}
\usepackage[scr=rsfs,cal=euler]{mathalfa}
\usepackage{xfrac}
\usepackage[unicode,hidelinks,bookmarks=false]{hyperref}
\newcommand{\ie}{i.e.\ }
\newcommand{\cf}{cf.\ }
\newcommand{\resp}{respectively\ }
\newcommand{\Subsection}{\subsection}
\providecommand{\nobreakdash}{\nobreak\hbox{-}}
\setlength{\emergencystretch}{2em}
\multlinegap0pt
\usepackage{bm}
\usepackage{mathtools}
\newtheorem{lemma}{Lemma}
\newtheorem{theorem}{Theorem}
\title{Correct mathematical models of joint filtration of two immiscible viscous liquids.}
\author{Anvarbek Meirmanov, Akbota Senkebayeva}
\date{Source: arXiv:2510.19502v6 (CC BY 4.0)}
\begin{document}
\maketitle

\noindent\emph{Source and licence.} Correct mathematical models of joint filtration of two immiscible viscous liquids.. Version arXiv:2510.19502v6 (CC BY 4.0), distributed by arXiv under the Creative Commons Attribution 4.0 International licence, http://creativecommons.org/licenses/by/4.0/, as recorded in the arXiv licence metadata for this exact version and shown on the abstract page https://arxiv.org/abs/2510.19502v6. Full licence notice: \texttt{licenses/arxiv-2510.19502-CC-BY-4.0.txt}.\par\smallskip

\noindent\textit{Editorial scope.} Counted unit: the article's lemma l5.3 (source0.tex lines 1160--1175). The article's complete original proof of it is delivered with it (source0.tex lines 1177--1395, one \texttt{proof} environment of 219 lines). No mathematics is written, repaired or re-derived: the statement and the proof are the article's own lines, in the article's own notation, and no retained line is substituted or re-flowed.  Retained uncounted as the source-local context the proof needs: lines 241--243; lines 662--669; lines 664--667; lines 916--920; lines 930--943; lines 955--957; lines 993--1003; lines 1119--1145.  Nothing was omitted from any retained span, so the package carries no omission record.  No retained line contains a citation, so the wrapper carries no bibliography and no bibliography entry.  No retained line contains a figure environment, an image-inclusion command or drawing code, so no image is delivered and the package declares no asset dependency.  Editorial formatting only, in addition: an AMS \texttt{amsart} preamble that restates the article's own declarations and macros used by the retained lines, namely the environments \texttt{lemma}, \texttt{theorem} on separate counters, together with the standard packages those lines need.  The retained environments print in this document as Lemma 1, Theorem 1 in order of appearance, whereas the article prints the same items with the numbering of its own full document.  The complete native source of the article as distributed in the arXiv e-print for this exact version is bundled unchanged as \texttt{sources/187576/source0.tex}.
\begin{equation}\label{eq2.1}
\mu^{\varepsilon}_{j}=\mu_{j}\chi^{\varepsilon}_{j},\,\, 0<\mu_{*}\leqslant\mu_{j}\leqslant\mu_{*}^{-1}<\infty,\,\,j=ol,\,sp,
\end{equation}

\begin{lemma}\label{l2.10}
For any $\widetilde{\boldsymbol{w}},\,\widetilde{\boldsymbol{v}}\in\mathbb{L}_{2}(\Omega)$ the H\"{o}lder's inequality  
\begin{equation}\label{eq2.9}
\|\widetilde{\boldsymbol{w}}\,\widetilde{\boldsymbol{v}}\|_{1,\Omega}\leqslant
\|\widetilde{\boldsymbol{w}}\|_{2,\Omega}\|\widetilde{\boldsymbol{v}}\|_{2,\Omega}
\end{equation}
holds.
\end{lemma}

\begin{equation}
\|\widetilde{\boldsymbol{w}}\,\widetilde{\boldsymbol{v}}\|_{1,\Omega}\leqslant
\|\widetilde{\boldsymbol{w}}\|_{2,\Omega}\|\widetilde{\boldsymbol{v}}\|_{2,\Omega}
\end{equation}

\begin{equation}\label{eq2.19}
\int_{0}^{t_{0}}\int_{\Omega}\big((\chi^{\varepsilon}_{sp}
\boldsymbol{v}^{\varepsilon}_{sp}+\chi^{\varepsilon}_{ol}
\boldsymbol{v}^{\varepsilon}_{ol})\cdot\boldsymbol{\varphi}\big)dxdt=0
\end{equation}

\begin{equation}\label{eq2.20}
\begin{aligned}
& 
\varepsilon^{2}\int_{0}^{t_{0}}\int_{\Omega}\big(\chi^{\varepsilon}_{sp}\mu_{sp}
\mathbb{D}(x,\boldsymbol{w}^{\varepsilon}_{sp})+\chi^{\varepsilon}_{ol}\mu_{ol}
\mathbb{D}(x,\boldsymbol{w}^{\varepsilon}_{ol})\big):
\mathbb{D}(x,\frac{\partial\boldsymbol{\varphi}}{\partial{t}})dxdt=
\\
& -\int_{0}^{t_{0}}\int_{\Omega}\Big(\big(\nabla({p}^{0}t)+\varepsilon^{2}
\chi^{\varepsilon}_{sp}\boldsymbol{w}^{\varepsilon}_{sp}+\chi^{\varepsilon}_{ol}
\varepsilon^{2}\boldsymbol{w}^{\varepsilon}_{ol}\big)
\cdot\frac{\partial\boldsymbol{\varphi}}{\partial{t}}\Big)dxdt,
\end{aligned}
\end{equation}

\begin{equation}\label{eq2.22}
\frac{\partial}{\partial{t}}(\rho{u})+\nabla\cdot(\rho{u}\boldsymbol{v})=0,
\end{equation}

\begin{theorem}\label{t3.1}
Let $\displaystyle\,0<r_{0}<\frac{1}{2}$, \, ${p}^{0}\in\mathbb{H}^{2+\alpha}(\overline{\Omega})$ and
$\displaystyle\,\frac{\partial{p^{0}}}{\partial{x}_{3}}(\boldsymbol{x})\,\leqslant\,0$
for $\boldsymbol{x}\in\Omega$.

Then under conditions \eqref{eq2.1} the problem $\mathbb{B}^{\varepsilon}$ has a unique weak solution
\begin{equation}\label{eq3.1}
\boldsymbol{w}^{\varepsilon}_{j}\in\mathbb{W}^{1,0}_{2}(\Omega_{j,T}),\,\boldsymbol{v}^{\varepsilon}_{j}
\in\mathbb{W}^{1,0}_{2}(\Omega_{j,T}),\,j=sp,\,ol.
\end{equation}
\end{theorem}

\begin{lemma}\label{l5.2}
Under the conditions of Theorem~\ref{t3.1} there exist unique functions $\boldsymbol{w}_{ol}$, $\boldsymbol{w}_{sp}$, $p_{ol}$, $p_{sp}$, $\displaystyle\,{\pi}_{ol}=\int_{0}^{t}p_{ol}(\boldsymbol{x},\tau)d\tau$ and $\displaystyle\,{\pi}_{sp}\int_{0}^{t}p_{sp}(\boldsymbol{x},\tau)d\tau$ such that
$p_{ol}, p_{sp},\boldsymbol{w}_{ol}, \boldsymbol{w}_{sp}\in\mathbb{W}^{1,0}_{2}(\Omega_{T})$,
$\boldsymbol{W}_{ol}, \boldsymbol{W}_{sp}\in\mathbb{L}_{2}\big(0,T;\mathbb{W}^{1}_{2}(\textbf{Y}_{f})\big)$.

1) The sequence $\{\widetilde{\boldsymbol{w}}^{\,\varepsilon}_{j}\}$ converges weakly to the function $\boldsymbol{w}_{j}$ and two-scale to the function
$\boldsymbol{W}_{j}(\boldsymbol{y};\boldsymbol{x},t),\,j=ol,\,sp$.

2) The sequences $\{\varepsilon\nabla\cdot\widetilde{\boldsymbol{w}}^{\,\varepsilon}_{j}\}$ converge two-scale to the functions $\nabla_{y}\cdot\boldsymbol{W}_{j},\,j=ol,\,sp$.

3) The sequences $\{\varepsilon\mathbb{D}(x,\widetilde{\boldsymbol{w}}^{\,\varepsilon}_{j})\}$ converge two-scale to the functions $\mathbb{D}(y,\boldsymbol{W}_{j}),\,j=ol,\,sp$.

4) The following a priori estimates hold true
\begin{multline}\label{eq5.1}
\|{\boldsymbol{w}}_{j}\|_{2,\Omega_{T}}+\|\boldsymbol{W}_{j}\|_{2,\textbf{Y}_{f}\times\Omega_{T}}+
\|\nabla_{y}\boldsymbol{W}_{j}\|_{2,{\textbf{Y}}_{f}\times\Omega_{T}}+
\|\mathbb{D}\big(y,\boldsymbol{W}_{j})\|_{2,{Y}_{f}\times\Omega_{T}}\leqslant
\\
\,{M}\sqrt{T}\|\nabla{p}^{0}\|_{2,\Omega},\,\,j=ol,\,sp,
\end{multline}
\begin{equation}\label{eq5.2}
\|{p}_{j}\|_{2,\Omega_{T}}+\|\nabla{p}_{j}\|_{2,\Omega_{T}}+\|{\pi}_{j}\|_{2,\Omega_{T}}+
\Big\|\nabla\frac{\partial{\pi}_{j}}{\partial{t}}\Big\|_{2,\Omega_{T}}\leqslant\,
{M}\sqrt{T}\|\nabla{p}^{0}\|_{2,\Omega},\,\,j=ol,\,sp,
\end{equation}
where the constant $M$ is bounded for bounded domains $\Omega$.
\end{lemma}

\begin{lemma}\label{l5.3}
Under the conditions of Theorem~\ref{t3.1} the limiting procedure in the integral identities \eqref{eq2.19} results in the following dynamic problem $\mathbb{H}$ for displacements and pressures of the liquid components consisting of Darcy's law of filtration
\begin{equation}\label{eq5.3}
\boldsymbol{w}_{j}=\int_{\textbf{Y}_{f}}\boldsymbol{W}_{j}dy=
-\frac{1}{\mu_{j}}(B)<\nabla_{x}({\pi}_{j}-{p}^{0}t)>,\,\,
\nabla_{x}\cdot\boldsymbol{w}_{j}=0,\,\,j=ol,\,sp
\end{equation}
for the liquid displacements $\boldsymbol{w}_{j}$ and the antiderivative ${\pi}_{j}$ of the pressure ${p}_{j},\,\,j=ol,\,sp$ in the domain $\Omega_{T}$, completed with the boundary conditions
\begin{equation}\label{eq5.4}
{\pi}_{j}(\boldsymbol{x},t)={p}^{0}t,\,\,j=ol,\,sp,\,\,\boldsymbol{x}\in{S}^{1}\cup{S}^{2},\,\,
\boldsymbol{w}_{j}\cdot\boldsymbol{n}=0,\,\,j=ol,\,sp,\,\,\boldsymbol{x}\in{S}^{0},\,\,0<t<T,
\end{equation}
where $\boldsymbol{n}$ is a normal vector to the boundary ${S}^{0}$.

The symmetric, strictly positive-definite constant matrix $(B)$ is given by the formula \eqref{eq5.16}.
\end{lemma}

\begin{proof}
First, we derive the continuity equations for the functions $\boldsymbol{w}_{j}$ and $\boldsymbol{W}_{j},\,j=ol,\,sp.$

To do this, let us consider the second integral identity in \eqref{eq2.22}
\begin{equation*}
0=\lim_{\varepsilon\rightarrow{0}}\int_{0}^{{t}_{0}}
\int_{\Omega}{\eta}(\nabla\cdot\boldsymbol{w}^{\varepsilon}_{j})dxdt=
-\int_{0}^{{t}_{0}}\int_{\Omega}(\nabla{\eta}\cdot\boldsymbol{w}_{j})dxdt,\,\,\,j=ol,\,sp.
\end{equation*}
This identity gives us continuity equations and boundary conditions for the displacements of oil and suspension
\begin{equation}\label{eq5.5}
\nabla\cdot\boldsymbol{w}_{j}=0,\,\,\boldsymbol{x}\in\Omega,\,\,(\boldsymbol{w}_{j}\cdot\boldsymbol{n})=0, \,\,\boldsymbol{x}\in{S}^{0},\,\,j=ol,\,sp.
\end{equation}
To derive the continuity equation for the unknown functions
$\boldsymbol{W}_{j}\big(\boldsymbol{y};\boldsymbol{x},t)\big),\,\,j=ol,\,sp$
(liquid displacements in the variables $\boldsymbol{y}$) we consider the integral identity \eqref{eq2.22} with arbitrary test functions $\displaystyle\,{\xi}=\varepsilon{\eta}(\boldsymbol{x},t)\phi(\frac{\boldsymbol{x}}{\varepsilon})$, where ${\eta}(\boldsymbol{x},t)$ is the same as before and $\phi(\boldsymbol{y})$ is 1-periodic in $\boldsymbol{y}$ function.

Furthermore, using statements~1) and~2) from the conditions of Lemma~\ref{l5.2} we obtain that
\begin{multline}\label{eq5.6}
0=\lim_{\varepsilon\rightarrow{0}}\int_{0}^{{t}_{0}}
\int_{\Omega}\varepsilon{\eta}\phi\chi^{\varepsilon}
(\nabla\cdot\widetilde{\boldsymbol{w}}_{j}^{\varepsilon})dxdt=
\\
\int_{0}^{{t}_{0}}\int_{\Omega}{\eta}\int_{\textbf{Y}_{f}}
\big(\phi\nabla_{y}\cdot\boldsymbol{W}_{j}\big)dydxdt,\,\,j=ol,\,sp.
\end{multline}
Due to the arbitrary choice of the functions ${\eta}$ and $\phi$ the last relation means that hold true the
continuity equation
\begin{equation}\label{eq5.7}
\nabla_{y}\cdot\boldsymbol{W}_{j}(\boldsymbol{y};\boldsymbol{x},t)=0,\,\,\,
(\boldsymbol{y};\boldsymbol{x},t)\in\textbf{Y}_{f}\times\Omega_{t_{0}},\,\,j=ol,\,sp
\end{equation}
together with boundary and normalization conditions
\begin{equation}\label{eq5.8}
\boldsymbol{W}_{j}=0,\,\,(\boldsymbol{y};\boldsymbol{x},t)\in\gamma\times\Omega_{t_{0}},\,\,j=ol,\,sp.
\end{equation}
Let additionally in \eqref{eq2.20} $\boldsymbol{\varphi}=0$ in $\Omega_{T}$, at $t=0$ and $t=t_{0}$,

$\displaystyle\,\frac{\partial\boldsymbol{\varphi}}{\partial{t}}=
\eta(\boldsymbol{x},t)\boldsymbol{\psi}(\frac{\boldsymbol{x}}{\varepsilon})$, where
$\eta\in\mathbb{C}^{\infty}(\overline{\Omega}_{T})$, $\eta(\boldsymbol{x},t)=0$ for $\boldsymbol{x}\in{S}^{0}$, $0<t<T$, $\boldsymbol{\psi}\in{W}^{1}_{2}(Y_{f})$,
\emph{supp}$\boldsymbol{\psi}\subset{Y}_{f}$, $\nabla_{y}\cdot\boldsymbol{\psi}=0$.

Then
\begin{equation*}
\begin{aligned}
& 
\mathbb{D}(x,{\eta}\frac{\partial\boldsymbol{\varphi}}{\partial{t}})=
\mathbb{D}(x,{\eta}\boldsymbol{\psi})=\frac{1}{2}\sum_{i,j=1}^{3}\Big({d}_{ij}\big(x,{\eta}
\boldsymbol{\psi}(\boldsymbol{y})\big)\boldsymbol{e}^{i}\otimes\boldsymbol{e}^{j}+
{d}_{ji}(x,{\eta}\boldsymbol{\psi}(\boldsymbol{y})\big)\boldsymbol{e}^{j}\otimes\boldsymbol{e}^{i}\Big)=
\\
& \frac{1}{2\varepsilon}{\eta}\big(\frac{\partial{\psi}_{j}}{\partial{y}_{i}}(\boldsymbol{y})
\boldsymbol{e}^{i}\otimes\boldsymbol{e}^{j}+\frac{\partial{\psi}_{i}}{\partial{y}_{j}}(\boldsymbol{y})\big)
\boldsymbol{e}^{j}\otimes\boldsymbol{e}^{i}\big)+
\frac{1}{2}\big(\frac{\partial{\eta}}{\partial{x}_{i}}{\psi}_{j}(\boldsymbol{y})
\boldsymbol{e}^{i}\otimes\boldsymbol{e}^{j}+
\frac{\partial{\eta}}{\partial{x}_{j}}{\psi}_{i}(\boldsymbol{y})
\boldsymbol{e}^{j}\otimes\boldsymbol{e}^{i}\big),
\\
& \varepsilon^{2}\mathbb{D}(x,\frac{\partial\boldsymbol{\varphi}}{\partial{t}})=
\varepsilon^{2}\mathbb{D}(x,{\eta}\boldsymbol{\psi})=
\eta\varepsilon\mathbb{D}\big(y,\boldsymbol{\psi}(\frac{\boldsymbol{x}}{\varepsilon})\big)+
\frac{\varepsilon^{2}}{2}(\nabla{\eta}\otimes\boldsymbol{\psi}+
\boldsymbol{\psi}\otimes\nabla{\eta}),
\\
& \nabla\cdot(\eta\boldsymbol{\psi})=
(\nabla{\eta}\cdot\boldsymbol{\psi})=-\frac{1}{\varepsilon}\eta(\nabla_{y}\cdot\boldsymbol{\psi})=0.
\end{aligned}
\end{equation*}
Next we consider functions
\begin{equation}\label{eq5.9}
\boldsymbol{A}_{j}(\boldsymbol{y};\boldsymbol{x},t)=
\nabla\cdot\big(\mu_{j}\mathbb{D}(y,\boldsymbol{W}_{j})\big)\,
\mbox{and}\,\boldsymbol{B}_{j}(\boldsymbol{x},t)=
(t\pi_{j}-{p}^{0})\int_{\textbf{Y}_{f}}\boldsymbol{\psi}dy,\,j=sp,\,ol
\end{equation}
and functional
\begin{equation*}
I^{\varepsilon}_{j}(\eta\boldsymbol{\psi})=\int_{0}^{t_{0}}\int_{\Omega}\chi^{\varepsilon}\eta
\int_{\textbf{Y}_{f}}\Big(\big(\varepsilon^{2}
\widetilde{\boldsymbol{w}}_{f}^{\varepsilon})\cdot\boldsymbol{\psi}\big)+
\big(\mu_{j}\varepsilon^{2}\mathbb{D}(x,\widetilde{\boldsymbol{w}}_{j}^{\varepsilon})+
(t\widetilde{{\pi}}^{\varepsilon}_{j}-{p}^{0})\mathbb{I}\big):\mathbb{D}(x,\boldsymbol{\psi})\Big)dydxdt.
\end{equation*}
In accordance with Lemma~5.2 and the integral identity \eqref{eq2.20} we get
\begin{equation}\label{eq5.10}
\begin{aligned}
& 
0={I}^{0}_{f}(\eta\boldsymbol{\psi})=
\lim_{\varepsilon\rightarrow{0}}I^{\varepsilon}_{f}(\eta\boldsymbol{\psi})=
-\lim_{\varepsilon\rightarrow{0}}\int_{0}^{t_{0}}\int_{\Omega}\chi^{\varepsilon}
\Big((\varepsilon^{2}\widetilde{\boldsymbol{w}}_{f}^{\varepsilon}+\nabla{p}^{0}t\cdot
\frac{\partial\boldsymbol{\varphi}}{\partial{t}})+
\\
& \big(\mu_{j}\varepsilon^{2}\mathbb{D}(x,\widetilde{\boldsymbol{w}}_{j}^{\varepsilon})-
(t\widetilde{{\pi}}^{\varepsilon}-{p}^{0})\mathbb{I}\big):
\mathbb{D}(x,\frac{\partial\boldsymbol{\varphi}}{\partial{t}})\Big)dxdt=
\\
& -\lim_{\varepsilon\rightarrow{0}}\int_{0}^{t_{0}}\int_{\Omega}\eta\chi^{\varepsilon}
\Big(\big(\nabla(\varepsilon\widetilde{\boldsymbol{w}}_{f}^{\varepsilon})\cdot\boldsymbol{\psi}\big)+
\mu_{j}\varepsilon\mathbb{D}(x,\widetilde{\boldsymbol{w}}_{j}^{\varepsilon})\big):
\mathbb{D}\big(y,\boldsymbol{\psi}(\frac{\boldsymbol{x}}{\varepsilon})\big)-
(t\widetilde{{\pi}}^{\varepsilon}-{p}^{0})\big(\nabla\eta\cdot
\boldsymbol{\psi}(\frac{\boldsymbol{x}}{\varepsilon})\big)\Big)dxdt=
\\
& -\int_{0}^{t_{0}}\int_{\Omega}\Big(\eta\big(\int_{Y_{f}}\big(-\nabla_{y}\cdot
\big(\mu_{j}\mathbb{D}(y,\boldsymbol{W}_{j})
\big)dy+(t{\pi}-{p}^{0})(\int_{\textbf{Y}_{f}}\boldsymbol{\psi}dy\cdot\nabla\eta)\Big)dxdt=
\\
& \int_{0}^{t_{0}}\int_{\Omega}\big((B)_{j}<\nabla\eta>-A_{j}\eta\big)dxdt=0.
\end{aligned}
\end{equation}
The last identity in \eqref{eq5.10}
\begin{equation}\label{eq5.11}
\int_{0}^{t_{0}}\int_{\Omega}\big((B)_{j}<\nabla\eta>-A_{j}\eta\big)dxdt=0
\end{equation}
means that functions ${\pi}_{j}\in\mathbb{W}^{1,0}_{2}(\Omega_{T})$ and identity
\eqref{eq5.11} takes the form of the differential equation
\begin{equation}\label{eq5.12}
\nabla_{y}\cdot\big(\mu_{j}\mathbb{D}(y,\boldsymbol{W}_{j})\big)=
-\nabla_{x}(t\pi_{j}-{p}^{0})(\boldsymbol{x},t)=
-\sum_{i=1}^{3}\frac{\partial{}}{\partial{x_{i}}}({\pi}_{j}-{p}^{0}t)(\boldsymbol{x},t)\boldsymbol{e}^{i},
\end{equation}
completed with the continuity equation \eqref{eq5.7}, boundary condition \eqref{eq5.8} and boundary condition \eqref{eq5.5}
\begin{equation}\label{eq5.13}
t\pi_{j}(\boldsymbol{x},t)-{p}^{0}=0,\,\,\,\boldsymbol{x}\in{S}^{1}\cup{S}^{2},\,\,0<t<T,
\end{equation}
which is a consequence of the identity \eqref{eq5.14}.

To solve the periodic boundary value problem \eqref{eq5.7}, \eqref{eq5.8}, \eqref{eq5.12} we use decomposition
\begin{equation}\label{eq5.14}
\boldsymbol{w}_{j}=\int_{\textbf{Y}_{f}}\boldsymbol{W}_{f}\big(\boldsymbol{y};\boldsymbol{x},t\big)dy=
-\frac{1}{\mu_{j}}\sum_{i=1}^{3}\int_{\textbf{Y}_{f}}\boldsymbol{W}_{f}^{(i)}(\boldsymbol{y})dy
\frac{\partial}{\partial{x_{i}}}({t\pi}_{j}-{p}^{0})
(\boldsymbol{x},t),
\end{equation}
where
\begin{equation}\label{eq5.15}
\begin{aligned}
& 
-\nabla\cdot\big(\mu_{j}\mathbb{D}(y,\boldsymbol{W}_{f}^{(i)})\big)=\boldsymbol{e}^{i},\,\,
\nabla\cdot\boldsymbol{W}_{f}^{(i)}=0,\,\,i=1,2,3,\,\,\boldsymbol{y}\in\textbf{Y}_{f},
\\
& \boldsymbol{W}_{f}^{(i)}=0,\,\,i=1,2,3,\,\,\boldsymbol{y}\in\gamma.
\end{aligned}
\end{equation}
The proof of the existence of solutions to the problem \eqref{eq5.15} is standard and follows from energy estimates
\begin{equation*}
\boldsymbol{W}_{f}^{(i)}\in\mathbb{W}^{(1,0}_{2}(\textbf{Y}_{f}),\,\,
\int_{\textbf{Y}_{f}}(|\boldsymbol{W}_{f}^{(i)}|^{2}+|\mathbb{D}(y,\boldsymbol{W}_{f}^{(i)})|^{2})dy
\leqslant\,{M},\,\,i=1,2,3,
\end{equation*}
which are the result of multiplying the equation in \eqref{eq5.15} by $\boldsymbol{W}_{f}^{(i)}$, integration by parts and use of the boundary condition in \eqref{eq5.15} and estimate \eqref{eq2.9} (see Lemma~\ref{l2.10}).

Next, we define the constant matrix $(B)$ as
\begin{equation}\label{eq5.16}
\begin{aligned}
& 
(B)=\sum_{i,j=1}^{3}({b}^{(i,j)}\boldsymbol{e}^{i}\otimes\boldsymbol{e}^{j}+
{b}^{(i,j)}\boldsymbol{e}^{j}\otimes\boldsymbol{e}^{i}),\,\,
\mathbb{B}<\boldsymbol{e}^{i},\boldsymbol{e}^{j}>={b}^{(i,j)},
\\
& {b}^{(i,j)}=
\int_{\textbf{Y}_{f}}(\boldsymbol{W}_{f}^{(i)}(\boldsymbol{y})\cdot\boldsymbol{e}^{j})dy.
\end{aligned}
\end{equation}
The matrix $(B)$ is obviously symmetric and strictly positive definite.
In fact, multiplying the equation in \eqref{eq5.16} by $\boldsymbol{W}_{f}^{(j)}$ gives
\begin{equation*}
\mu_{j}\int_{{Y}_{f}}\mathbb{D}(y,\boldsymbol{W}_{f}^{(i)}):
\mathbb{D}(y,\boldsymbol{W}_{f}^{(j)})dy=
\int_{\textbf{Y}_{f}}(\boldsymbol{W}_{f}^{(j)}\cdot\boldsymbol{e}^{i})dy.
\end{equation*}
Then the equality
\begin{equation*}
\mu_{j}\int_{\textbf{Y}_{f}}\mathbb{D}(y,\boldsymbol{W}_{f}^{(i)}):
\mathbb{D}(y,\boldsymbol{W}_{f}^{(j)})dy=
\int_{\textbf{Y}_{f}}\mu_{j}\mathbb{D}(y,\boldsymbol{W}_{f}^{(j)}):
\mathbb{D}(y,\boldsymbol{W}_{f}^{(i)})dy
\end{equation*}
implies the equality
\begin{equation}\label{eq5.17}
\int_{\textbf{Y}_{f}}(\boldsymbol{W}_{f}^{(i)}\cdot\boldsymbol{e}^{j})dy=
\int_{\textbf{Y}_{f}}(\boldsymbol{W}_{f}^{(j)}\cdot\boldsymbol{e}^{i})dy,
\end{equation}
which means the symmetry of the matrix $(B)_{f}$.

To prove the strict positive definiteness of the matrix $(B)$ we put
\begin{equation*}
\boldsymbol{W}_{f}(\boldsymbol{\xi})=\sum_{i=1}^{3}\boldsymbol{W}_{f}^{(i)}\xi_{i}
\end{equation*}
for any vector $\boldsymbol{\xi}\in\mathbb{R}^{3}$ and any function
$\boldsymbol{\varphi}\in\stackrel{\!\!\circ}{\mathbb{W}}^{1,0}_{2}({Y}_{f})$ and consider
the integral identity
\begin{equation}\label{eq5.18}
\int_{\textbf{Y}_{f}}\mathbb{D}\big(y,\boldsymbol{W}_{f}(\boldsymbol{\xi})\big):
\mathbb{D}(y,\boldsymbol{\varphi})dy=
\int_{\textbf{Y}_{f}}(\boldsymbol{\varphi}\cdot\boldsymbol{\xi})dy.
\end{equation}
Then for $\boldsymbol{\varphi}=\boldsymbol{W}_{f}(\boldsymbol{\xi})$ one has
\begin{equation}\label{eq5.19}
\begin{aligned}
& 
\int_{\textbf{Y}_{f}}|\mathbb{D}(y,\boldsymbol{W}_{f}(\boldsymbol{\xi})|^{2}dy=
\int_{\textbf{Y}_{f}}\big(\boldsymbol{\xi}\cdot\boldsymbol{W}_{f}(\boldsymbol{\xi})\big)dy=
\\
& \sum_{i,j=1}^{3}{\xi}_{i}{\xi}_{j}\int_{\textbf{Y}_{f}}(\boldsymbol{W}_{f}^{(i)}
\cdot\boldsymbol{e}^{j})dy=\sum_{i,j=1}^{3}{b}_{i,j}{\xi}_{i}{\xi}_{j}>0.
\end{aligned}
\end{equation}
It is evident that the equality $\displaystyle\,\sum_{i,j=1}^{3}{b}_{i,j}{\xi}_{i}{\xi}_{j}=0$
implies the equalities $\boldsymbol{W}_{f}=0$ and $\mathbb{D}\big(y,\boldsymbol{W}_{f}(\boldsymbol{\xi})\big)=0$ in $\textbf{Y}_{f}$, which is impossible.

Therefore, we can limit ourselves to the case $|\boldsymbol{\xi}|=1$. This fact immediately leads to the inequality
\begin{equation*}
\sum_{i,j=1}^{3}{b}_{i,j}{\xi}_{i}{\xi}_{j}\geqslant\,\alpha_{0}=\mbox{const}>0.
\end{equation*}
\end{proof}

\end{document}
